\section{矢量分析习题解答}

\begin{enumerate}[label=\arabic*.,itemsep=10pt,topsep=5pt,parsep=2pt]

% --------------- 第1题（作业题号4）---------------
\item 设有标量场 $\varphi = 2xy^2 - z^3$，求其在点 $M(2,-1,1)$ 沿 $\overrightarrow{MP}\big(P(3,1,-1)\big)$ 方向的方向导数，以及该点方向导数的最大值。

\textbf{解答：}
\begin{enumerate}[label=(\arabic*),itemsep=5pt,topsep=2pt,parsep=2pt]
\item 求取梯度
\[
\nabla\varphi = \frac{\partial \varphi}{\partial x}\vec{e}_x+\frac{\partial \varphi}{\partial y}\vec{e}_y+\frac{\partial \varphi}{\partial z}\vec{e}_z
= 2y^2\vec{e}_x + 4xy\vec{e}_y - 3z^2\vec{e}_z
\]
将 $M(2,-1,1)$ 代入坐标得：
\[
\nabla\varphi\big|_M = 2\vec{e}_x - 8\vec{e}_y - 3\vec{e}_z
\]

\item 方向矢量 $\overrightarrow{MP}=(1,2,-2)$，其模长
\[
|\overrightarrow{MP}|=\sqrt{1^2+2^2+(-2)^2}=3
\]
单位矢量为
\[
\vec{e}_l = \frac{1}{3}\vec{e}_x + \frac{2}{3}\vec{e}_y - \frac{2}{3}\vec{e}_z
\]

\item 由方向导数公式 $\displaystyle\frac{\partial \varphi}{\partial l}=\nabla\varphi\cdot\vec{e}_l$
\[
\left.\frac{\partial\varphi}{\partial l}\right|_M
=2\cdot\frac13+(-8)\cdot\frac23+(-3)\cdot\left(-\frac23\right)
=-\frac{8}{3}
\]

\item 方向导数最大值出现在梯度同向，大小为梯度模长
\[
\left|\nabla\varphi\big|_M\right|
=\sqrt{2^2+(-8)^2+(-3)^2}
=\sqrt{77}
\]
\end{enumerate}
\textbf{结论：}方向导数$\boxed{-\dfrac{8}{3}}$，最大值$\boxed{\sqrt{77}}$。

% --------------- 第2题（作业题号5）---------------
\item 求标量场 $\varphi = x^3y^4z^2$ 梯度场的散度。

\textbf{解答：}
先求梯度：
\[
\nabla\varphi 
=\frac{\partial\varphi}{\partial x}\vec{e}_x+\frac{\partial\varphi}{\partial y}\vec{e}_y+\frac{\partial\varphi}{\partial z}\vec{e}_z
= 3x^2y^4z^2\vec{e}_x + 4x^3y^3z^2\vec{e}_y + 2x^3y^4z\vec{e}_z
\]
再求梯度的散度（拉普拉斯算子）：
\[
\begin{aligned}
\nabla\cdot(\nabla\varphi)
&=\frac{\partial}{\partial x}(3x^2y^4z^2)+\frac{\partial}{\partial y}(4x^3y^3z^2)+\frac{\partial}{\partial z}(2x^3y^4z) \\
&= 6xy^4z^2 + 12x^3y^2z^2 + 2x^3y^4
\end{aligned}
\]
\textbf{结论：}$\boxed{6xy^4z^2 + 12x^3y^2z^2 + 2x^3y^4}$。

% --------------- 第3题（作业题号7）---------------
\item 设矢量场 $\vec{A} = x^2\vec{e}_x + y^3\vec{e}_y + (3z - x)\vec{e}_z$，求：(1) $M_1(1,0,-1)$ 处散度；(2) $M_2(1,-1,-1)$ 处旋度。

\textbf{解答：}
\begin{enumerate}[label=(\arabic*),itemsep=5pt,topsep=2pt,parsep=2pt]
\item 散度计算
\[
\nabla\cdot\vec{A} 
=\frac{\partial A_x}{\partial x}+\frac{\partial A_y}{\partial y}+\frac{\partial A_z}{\partial z}
=\frac{\partial(x^2)}{\partial x}+\frac{\partial(y^3)}{\partial y}+\frac{\partial(3z-x)}{\partial z}
=2x+3y^2+3
\]
代入 $M_1(1,0,-1)$：
\[
\nabla\cdot\vec{A}\big|_{M_1}=2\times 1+3\times 0^2+3=5
\]

\item 旋度计算
\[
\nabla\times\vec{A} =
\begin{vmatrix}
\vec{e}_x & \vec{e}_y & \vec{e}_z \\
\dfrac{\partial}{\partial x} & \dfrac{\partial}{\partial y} & \dfrac{\partial}{\partial z} \\
x^2 & y^3 & 3z-x
\end{vmatrix}
=\vec{e}_y
\]
该旋度为常矢量，与坐标无关，故 $\nabla\times\vec{A}\big|_{M_2}=\vec{e}_y$。
\end{enumerate}
\textbf{结论：}散度$\boxed{5}$，旋度$\boxed{\vec{e}_y}$。

% --------------- 第4题（作业题号8）---------------
\item 已知电场强度 $\vec{E} = E_0\cos\theta\vec{e}_r - E_0\sin\theta\vec{e}_\theta$，求其散度与旋度。

\textbf{解答：}
在球坐标系中计算散度与旋度，代入化简可得：
\[
\nabla\cdot\vec{E} = 0,\quad \nabla\times\vec{E} = \vec{0}
\]
\textbf{结论：}$\boxed{\nabla\cdot\vec{E}=0,\ \nabla\times\vec{E}=\vec{0}}$。

% --------------- 第5题（作业题号9）---------------
\item 矢量场 $\vec{A}=(x^2-2xy)\vec{e}_x+(y^2-2yz)\vec{e}_y+z(z-2x+1)\vec{e}_z$，求其对半径为$a$、球心在原点的球面外侧的通量。

\textbf{解答：}
由高斯通量定理：
\[
\varPhi=\oiint_S \vec{A}\cdot d\vec{S}=\iiint_V \nabla\cdot\vec{A}\,\mathrm{d}V
\]
计算散度：
\[
\nabla\cdot\vec{A}=1
\]
因此：
\[
\varPhi=\iiint_V 1\,\mathrm{d}V=\frac{4}{3}\pi a^3
\]
\textbf{结论：}通量$\boxed{\dfrac{4}{3}\pi a^3}$。

% --------------- 第6题（作业题号10）---------------
\item 求矢量场 $\vec{A}=xyz(\vec{e}_x+\vec{e}_y+\vec{e}_z)$ 在点 $M(1,3,2)$ 处的旋度，以及绕方向 $\vec{e}_n=\frac{1}{3}(\vec{e}_x+2\vec{e}_y+2\vec{e}_z)$ 的环量面密度。

\textbf{解答：}
旋度计算并代入 $M(1,3,2)$：
\[
\nabla\times\vec{A}\big|_M
=-\vec{e}_x-3\vec{e}_y+4\vec{e}_z
\]
环量面密度为旋度在法向的投影：
\[
\mu_n=(\nabla\times\vec{A})\cdot\vec{e}_n
=\frac{1}{3}
\]
\textbf{结论：}旋度$\boxed{-\vec{e}_x-3\vec{e}_y+4\vec{e}_z}$，环量面密度$\boxed{\dfrac{1}{3}}$。

\end{enumerate}